% ============================================================
% 07resultados.tex
% Experimentos y resultados (Etapa 2)
% Fuente de las cifras: notebooks/06_evaluacion_final.ipynb y
% docx/03Evidence/04-11-26_ResultadosEvaluacionFinal.md
% ============================================================

\section{Experimentos y Resultados}
\label{sec:resultados}

Esta sección explica cómo se entrenaron y evaluaron los modelos, y presenta
los resultados sobre el conjunto de prueba. En todas las métricas cada
partida pesa lo mismo, sin importar cuántos estados tenga. Los intervalos
entre corchetes son intervalos de confianza del 95\,\%, calculados
remuestreando partidas 1000 veces.

\subsection{Ejecución del protocolo}
\label{subsec:ejecucion}

\noindent\textbf{Datos.} Las repeticiones de los diez días se convirtieron
en una tabla con una fila por estado, vista desde el jugador que actúa. La
tabla tiene 55\,196 partidas, 7\,286\,386 estados utilizables y las 62
características públicas de la Tabla~\ref{tab:diccionario}. Las etapas se
definieron con el entrenamiento: inicio hasta el turno 5, medio del 6 al 9
y final desde el 10. Cada partida quedó completa en un solo subconjunto, y
24 pruebas automáticas comprueban que no haya fuga
(Sección~\ref{subsec:leakage}).

\noindent\textbf{Modelos.} Se entrenaron cinco modelos, de menor a mayor
complejidad:
\begin{itemize}
  \item \textbf{B0:} predice siempre la proporción de victorias del
        entrenamiento.
  \item \textbf{B1:} convierte una diferencia fija de recursos en
        probabilidad.
  \item \textbf{B2:} regresión logística regularizada ($C = 0.1$).
  \item \textbf{B3:} bosque aleatorio.
  \item \textbf{B4:} potenciación del gradiente
        (\texttt{HistGradientBoostingClassifier}). Es el modelo principal.
\end{itemize}

\noindent\textbf{Búsqueda de hiperparámetros de B4.} Se usó Optuna con 20
ensayos y validación cruzada de tres pliegues por partida. La búsqueda usó
solo las 11\,285 partidas que son de entrenamiento en los tres escenarios,
y minimizó la pérdida logarítmica. Se buscaron los cuatro parámetros que
controlan la capacidad del modelo: tasa de aprendizaje entre 0.02 y 0.2,
regularización $L_2$ entre $10^{-6}$ y 10, número máximo de hojas en
$\{8, 16, 31\}$ y mínimo de muestras por hoja en $\{20, 50, 100\}$. El
número de árboles no se buscó: lo decide la parada temprana. Se limitó a 20
ensayos porque el cómputo disponible era una máquina de ocho núcleos sin
GPU. Con la mejor configuración, B4 se reentrenó con todo el entrenamiento
de cada escenario. Como control, LightGBM con el mismo presupuesto dio el
mismo AUC en validación (diferencia $\le 0.002$), así que el resultado no
depende de la librería.

\noindent\textbf{Evaluación.} El plan de evaluación se fijó antes de mirar
el conjunto de prueba. Antes de abrirlo se comprobó que pasaran 109 pruebas
automáticas, que los modelos reprodujeran exactamente sus métricas de
validación y que SHAP no mostrara fuga (Sección~\ref{subsec:xai}). El
conjunto de prueba se usó una sola vez.

\subsubsection{Cambios respecto del diseño}
\label{subsec:desviaciones}

\begin{itemize}
  \item Se usó una sola semilla. No se midió cuánto varían los resultados
        al repetir el entrenamiento.
  \item El escenario C reserva arquetipos de mazo, definidos por su
        Pok\'emon principal, y no mazos completos.
  \item Las ablaciones quitaron tres familias de características: progreso,
        recursos y tablero, y diferenciales. El contexto de mazo se evaluó
        en la prueba de H2. No se entrenó la variante de 64 características.
  \item La explicabilidad usó solo SHAP sobre B4. No se reportan los
        coeficientes de B2 ni la importancia por permutación.
\end{itemize}

\subsection{Desempeño frente a las referencias}
\label{subsec:resgeneral}

La Tabla~\ref{tab:resultados} muestra el AUC en prueba. El AUC mide si el
modelo asigna más probabilidad a los estados ganadores que a los
perdedores: 0.5 es azar y 1 es perfecto. Cada modelo supera al anterior en
los tres escenarios, y B4 es el mejor. Todas las diferencias entre B4 y
B0--B3 tienen intervalos que excluyen el cero, así que no se deben al azar
de la muestra. Frente a B2, B4 gana 0.071, 0.043 y 0.036 de AUC en A, B y
C.

\begin{table}[!t]
\centering
\caption{AUC en prueba, con intervalo del 95\,\%, y pérdidas de B4.}
\label{tab:resultados}
\footnotesize
\setlength{\tabcolsep}{3pt}
\begin{tabular}{@{}lccc@{}}
\toprule
\textbf{Modelo} & \textbf{A} & \textbf{B} & \textbf{C}\\
\midrule
B0 & 0.500 & 0.500 & 0.500\\
B1 & 0.683 {\scriptsize[0.674, 0.691]} & 0.718 {\scriptsize[0.711, 0.726]} & 0.705 {\scriptsize[0.697, 0.713]}\\
B2 & 0.762 {\scriptsize[0.754, 0.768]} & 0.746 {\scriptsize[0.740, 0.753]} & 0.751 {\scriptsize[0.745, 0.758]}\\
B3 & 0.804 {\scriptsize[0.798, 0.809]} & 0.771 {\scriptsize[0.765, 0.777]} & 0.767 {\scriptsize[0.761, 0.773]}\\
\textbf{B4} & \textbf{0.833} {\scriptsize[0.827, 0.838]} & \textbf{0.789} {\scriptsize[0.784, 0.795]} & \textbf{0.787} {\scriptsize[0.782, 0.793]}\\
\midrule
B4 log loss & 0.494 & 0.551 & 0.550\\
B4 Brier & 0.166 & 0.188 & 0.188\\
\bottomrule
\end{tabular}
\end{table}

En A y B el AUC de prueba es casi igual al de validación (0.835 y 0.786).
En C baja de 0.836 a 0.787. Es lo esperado: la validación de C tiene
arquetipos que el modelo ya conoce y la prueba tiene arquetipos nuevos. La
literatura reporta AUC de 0.78 a 0.80 (Sección~\ref{sec:baseline}). B4 llega
a esa banda desde el turno 7 y alcanza 0.937 en la etapa final de A
(Fig.~\ref{fig:aucturno}). Esta comparación es solo orientativa, porque los
juegos y los datos son distintos.

\subsection{Evolución con la partida (H1)}
\label{subsec:resh1}

\begin{figure}[!t]
\centering
\includegraphics[width=\columnwidth]{res_auc_por_turno.png}
\caption{AUC de B4 por turno en la prueba del escenario A, con banda del
95\,\%. La franja gris marca el rango de la literatura.}
\label{fig:aucturno}
\end{figure}

El AUC sube con cada turno: 0.58 en el turno 1 y cerca de 0.92 desde el
turno 15 (Fig.~\ref{fig:aucturno}). Las partidas largas podrían ser más
fáciles de predecir, y eso inflaría la etapa final. Para descartarlo, la
Tabla~\ref{tab:h1} también mide solo las partidas que llegan a las tres
etapas. En ellas el AUC sube unos 0.32 de inicio a final en los tres
escenarios. Además, en la etapa de inicio el AUC ya es claramente mayor que
0.5.

\begin{table}[!t]
\centering
\caption{AUC de B4 por etapa en prueba, y aumento de inicio a final en las
partidas que llegan a las tres etapas.}
\label{tab:h1}
\footnotesize
\begin{tabular}{@{}lccc@{}}
\toprule
& \textbf{A} & \textbf{B} & \textbf{C}\\
\midrule
Inicio & 0.684 & 0.630 & 0.623\\
Medio & 0.828 & 0.774 & 0.787\\
Final & 0.937 & 0.915 & 0.919\\
\midrule
Partidas en las tres etapas & 6\,839 & 8\,269 & 7\,412\\
Final $-$ inicio & 0.317 & 0.325 & 0.328\\
\bottomrule
\end{tabular}
\end{table}

\subsection{Identidad frente a estado (H2) y ablaciones}
\label{subsec:resh2}

Para probar H2 se entrenó un modelo restringido que solo conoce el turno,
quién empezó y qué Pok\'emon están activos. La Tabla~\ref{tab:h2} lo
compara con B4 en la etapa de inicio. En A, los dos modelos empatan. En B y
C, B4 es claramente mejor. En la etapa final de A, el modelo restringido
llega a 0.769 y B4 a 0.937.

\begin{table}[!t]
\centering
\caption{H2: AUC en la etapa de inicio, B4 frente al modelo restringido.}
\label{tab:h2}
\footnotesize
\setlength{\tabcolsep}{4pt}
\begin{tabular}{@{}lcccl@{}}
\toprule
& \textbf{B4} & \textbf{Restr.} & \textbf{Diferencia [IC]} & \textbf{H2}\\
\midrule
A & 0.684 & 0.682 & 0.002 {\scriptsize[$-$0.008, 0.011]} & Refutada\\
B & 0.630 & 0.589 & 0.041 {\scriptsize[0.031, 0.051]} & Sostenida\\
C & 0.623 & 0.535 & 0.088 {\scriptsize[0.077, 0.099]} & Sostenida\\
\bottomrule
\end{tabular}
\end{table}

Las ablaciones quitan una familia de características y miden cuánto cae el
AUC en validación. Sin recursos y tablero, el AUC cae entre 0.05 y 0.06: es
la familia esencial. Sin progreso, cae como mucho 0.003. Sin diferenciales
no cae, porque el modelo puede calcular esas diferencias por sí mismo. B4 no
se cambió con estos resultados, porque elegir características mirando la
validación sesgaría la evaluación.

\subsection{Calibración y recalibración}
\label{subsec:rescalibracion}

La calibración es el módulo adicional que eligió el equipo. Un modelo está
calibrado si sus probabilidades coinciden con lo que ocurre: de los estados
a los que asigna 0.8, debería ganar el 80\,\%. Se mide con el error de
calibración esperado (ECE), donde 0 es perfecto. Para corregir un modelo
mal calibrado se probaron dos métodos, Platt e isotónica, ajustados solo con
el conjunto de calibración. El método de cada modelo se eligió antes de
abrir la prueba.

\begin{figure}[!t]
\centering
\includegraphics[width=\columnwidth]{res_fiabilidad.png}
\caption{Diagramas de fiabilidad de B4 en prueba. La diagonal es la
calibración perfecta. Se muestran el modelo sin corrección y con las
correcciones isotónica y de Platt.}
\label{fig:fiabilidad}
\end{figure}

\begin{table}[!t]
\centering
\caption{Calibración en prueba, antes $\rightarrow$ después de corregir.}
\label{tab:calibracion}
\footnotesize
\setlength{\tabcolsep}{3pt}
\begin{tabular}{@{}lccc@{}}
\toprule
\textbf{Modelo} & \textbf{Log loss} & \textbf{Brier} & \textbf{ECE}\\
\midrule
A, B4 (ninguno) & 0.494 & 0.166 & 0.002\\
A, B3 (Platt) & 0.534 $\rightarrow$ 0.531 & 0.181 $\rightarrow$ 0.180 & 0.030 $\rightarrow$ 0.008\\
B, B4 (isotónica) & 0.551 $\rightarrow$ 0.547 & 0.188 $\rightarrow$ 0.187 & 0.037 $\rightarrow$ 0.018\\
B, B3 (Platt) & 0.564 $\rightarrow$ 0.568 & 0.193 $\rightarrow$ 0.195 & 0.013 $\rightarrow$ 0.034\\
\bottomrule
\end{tabular}
\end{table}

La Fig.~\ref{fig:fiabilidad} y la Tabla~\ref{tab:calibracion} muestran tres
resultados:
\begin{itemize}
  \item En A, B4 ya está calibrado (ECE 0.002) y no necesita corrección. El
        boosting suele distorsionar las probabilidades
        \cite{niculescumizil_probabilities_2005}, pero aquí no ocurre.
  \item En B, B4 es demasiado confiado: cuando predice 0.87, gana alrededor
        del 80\,\%. La corrección isotónica reduce el ECE a la mitad
        (0.037 a 0.018) sin cambiar el AUC.
  \item La misma estrategia empeora a B3 en B. Una corrección ajustada con
        el 11 de julio no siempre sirve para el 15 y el 22.
\end{itemize}

\subsection{Degradación bajo cambio temporal (H3)}
\label{subsec:resh3}

H3 predice que, cuando cambia la población, las probabilidades se
estropean más que el orden. Para probarlo se usaron las 1\,405 partidas que
son de prueba en A y en B (15 y 22 de julio). Sobre esas mismas partidas se
compara el modelo de A, que vio esos días, con el de B, que solo vio días
anteriores. La pérdida extra del modelo de B se divide en dos partes
(Fig.~\ref{fig:descomposicion}): la que se debe a peor calibración
($\Delta$MCB) y la que se debe a peor ordenamiento ($\Delta$DSC).

\begin{table}[!t]
\centering
\caption{H3: pérdida extra del modelo de B frente al de A en las mismas
partidas.}
\label{tab:h3}
\footnotesize
\setlength{\tabcolsep}{3pt}
\begin{tabular}{@{}lcc@{}}
\toprule
\textbf{Medida} & \textbf{Log loss} & \textbf{Brier}\\
\midrule
Calibración ($\Delta$MCB) & 0.008 {\scriptsize[0.005, 0.012]} & 0.002 {\scriptsize[0.002, 0.003]}\\
Ordenamiento ($\Delta$DSC) & 0.030 {\scriptsize[0.023, 0.036]} & 0.012 {\scriptsize[0.009, 0.014]}\\
Contraste & $-$0.021 {\scriptsize[$-$0.028, $-$0.015]} & $-$0.009 {\scriptsize[$-$0.012, $-$0.007]}\\
\bottomrule
\end{tabular}
\end{table}

La Tabla~\ref{tab:h3} muestra lo contrario de H3. La calibración empeora un
poco, pero el ordenamiento empeora entre 3.5 y 4.9 veces más; el AUC baja
0.027. Comparar porcentajes de caída habría dado la conclusión opuesta
(22.9\,\% en log loss contra 8.9\,\% en AUC), porque las dos métricas tienen
escalas distintas. Por eso el protocolo decide con la descomposición. El
modelo de B tampoco empeora día a día: su AUC oscila entre 0.775 y 0.792 del
7 al 22 de julio.

\subsection{Curvas de aprendizaje y sobreajuste}
\label{subsec:rescurvas}

B4 se entrenó con el 25, 50, 75 y 100\,\% de las partidas de
entrenamiento. El último 25\,\% sube el AUC de validación como mucho 0.004,
así que más datos aportarían poco. La diferencia entre el AUC de
entrenamiento y el de validación se mantiene en unos 0.034 en A y C a
medida que crecen los datos, así que no hay sobreajuste creciente. En B esa
diferencia es de 0.09 y tampoco cambia con el tamaño. No es sobreajuste: se
debe a que validación tiene días distintos a los de entrenamiento.

\subsection{Análisis de errores}
\label{subsec:reserrores}

Se clasificó cada estado como victoria si $p \ge 0.5$. Los principales
hallazgos son:
\begin{itemize}
  \item Al inicio, el modelo sobreestima la victoria de quien actúa: los
        falsos positivos (0.22--0.25 de los estados) superan a los falsos
        negativos (0.15--0.17). En la etapa final se igualan.
  \item Los errores muy seguros ($p \ge 0.9$ con derrota o $p \le 0.1$ con
        victoria) son menos del 1\,\% de los estados, y son más frecuentes
        en B y C que en A.
  \item Las predicciones muy seguras aciertan entre el 80 y el 90\,\% al
        inicio y entre el 96 y el 98\,\% al final. El peor caso es el inicio
        de B, con 79.7\,\%.
  \item Ninguna característica aislada explica los errores muy seguros:
        suelen ocurrir en estados poco avanzados de la partida.
\end{itemize}

\subsection{Explicabilidad}
\label{subsec:xai}

\begin{figure}[!t]
\centering
\includegraphics[width=0.92\columnwidth]{res_shap.png}
\caption{Valores SHAP del modelo final (B4 de A) en partidas de validación.
Cada punto es un estado. A la derecha del cero la característica aumenta la
probabilidad de ganar; el color indica si su valor es alto o bajo.}
\label{fig:shap}
\end{figure}

SHAP indica cuánto aporta cada característica a cada predicción. Se aplicó
al modelo final, B4 entrenado en A, sobre partidas de validación
(Fig.~\ref{fig:shap}). Recursos y tablero explican el 72.8\,\% de la
predicción, los diferenciales el 20.9\,\% y el progreso el 6.3\,\%. Las
características más importantes son la diferencia de premios (0.113), los
descartes propios y rivales (0.082 y 0.077), los premios propios y los
puntos de vida máximos de ambos activos. La diferencia de premios gana peso
con la partida: pasa de 0.057 al inicio a 0.149 al final.

SHAP también sirvió para buscar fuga antes de abrir la prueba. Si una
característica concentra demasiada importancia, suele ser porque filtra el
resultado. No se encontró fuga: ninguna columna prohibida participa, la
característica principal concentra 0.113 (el límite era 0.5) y la mejor
característica sola llega a un AUC de 0.690 (el límite era 0.85).